\documentclass{nbeamer}

\usepackage{nbeamer-code}
\title{nbeamer-code 使用手册}
\subtitle{代码排版}
\author{XIAN}
\institute{纳米树研究室 Namishu Research Lab}
\date{\today}

\begin{document}
\begin{frame}
  \titlepage
\end{frame}

\begin{frame}
  \title{目录}
  \usebeamertemplate{contents page}
\end{frame}

\section{使用}
\begin{frame}[fragile]{命令介绍}
行内引用用 \incode{\incode}，独立代码块用 \incode{ncode} 环境。在导言区加载 \incode{\usepackage{nbeamer-code}}，给展示代码的页面加上 \incode{[fragile]}。
\begin{ncode}[escapeinside={(*@}{@*)}]
\begin{frame}[fragile]{计算中点}
调用 \incode{midpoint(a, b)} 可以得到两个数的中点。
\begin{ncode}
def midpoint(a, b):
    return (a + b) / 2
(*@\incode{\end{ncode}}@*)
(*@\incode{\end{frame}}@*)
\end{ncode}
函数名随句子排在行内，函数定义单独成块。
\end{frame}

\begin{frame}[fragile]{行内代码}
要在正文中写出一条 LaTeX 命令，把它放进 \incode{\incode{...}}：
\begin{ncode}
\incode{\vspace{\nbgrid}}
\end{ncode}
用 \incode{\vspace{\nbgrid}} 在段间增加一格留白。

反斜杠和花括号都会原样显示。要引用单独的左花括号，可以用分隔符写成 \incode+\incode|{|+。
\end{frame}

\begin{frame}[fragile]{代码块}
\begin{columns}
\begin{column}{11\nbgrid}
\textbf{效果}
\begin{ncode}[language=Python]
def midpoint(a, b):
    total = a + b
    return total / 2
\end{ncode}
\end{column}%
\hspace{2\nbgrid}%
\begin{column}{15\nbgrid}
\textbf{源码}
\begin{ncode}[escapeinside={(*@}{@*)}]
\begin{ncode}[language=Python]
def midpoint(a, b):
    total = a + b
    return total / 2
(*@\incode{\end{ncode}}@*)
\end{ncode}
按原有缩进写入代码。可选参数 \incode{language=Python} 启用语法高亮。左栏的 \incode{def} 和 \incode{return} 因此显示为粗体。
\end{column}
\end{columns}
\end{frame}

\section{排版}
\begin{frame}[fragile]{自动换行}
\begin{columns}
\begin{column}{13\nbgrid}
\begin{ncode}[language=Python]
def area(width, height, scale):
    return width * height * scale * scale
\end{ncode}
\end{column}%
\hspace{2\nbgrid}%
\begin{column}{13\nbgrid}
左栏的 \incode{return} 语句在源码中只有一行。由于这行代码超过栏宽，\incode{ncode} 会自动换行。

源码中无需手动断行。排版后，相邻两行的基线仍相距一格。
\end{column}
\end{columns}
\end{frame}

\end{document}
